%% kapa: exam=1 round=30 year=2019 period=2 form=A subject=appraisal_law subject_ko=감정평가관계법규 seq=18 number=18
%% source: Q-Net 감정평가사 기출문제 (cst00309) artlSeq=5205205 file=extracted/1008/5205205/2019년 제30회 감정평가사 1차시험 문제지/2019년 제30회 감정평가사 1차시험 문제지/2019년도 제30회 감정평가사 1차 2교시 A형.hwp
%% conversion: kordoc 4.24.51 -> tools/gen1.py (auto); review=proofread-sonnet; flags=none
\begin{question}{18}
부동산 가격공시에 관한 법령상 공동주택가격의 조사ㆍ산정 및 공시에 관한 설명으로 옳지 않은 것은?
\begin{choices}
  \item 국토교통부장관이 공동주택가격을 조사ㆍ산정하고자 할 때에는 한국감정원에 의뢰한다.
  \item 국토교통부장관은 3월 31일에 대지가 합병된 공동주택의 공동주택가격을 그 해 6월 1일까지 산정ㆍ공시하여야 한다.
  \item 공시기준일 이후「건축법」에 따른 용도변경이 된 공동주택은 공동주택가격의 공시기준일을 다르게 할 수 있는 공동주택에 해당한다.
  \item 공동주택가격의 공시에는 공동주택의 면적이 포함되어야 한다.
  \item 아파트에 해당하는 공동주택은 국세청장이 국토교통부장관과 협의하여 그 공동주택가격을 별도로 결정ㆍ고시할 수 있다.
\end{choices}
\end{question}
\answer{2}
